\chapter*{A Note from the Author}
\addcontentsline{toc}{chapter}{A Note from the Author}
\markboth{A Note from the Author}{A Note from the Author}

I started this a very long time ago. As a boy in primary school I wrote a letter
to Dr.\ John Renner Hansen at the Niels Bohr Institute in Copenhagen, pressing on
him my formulas for cloud-chamber measurements and asking whether he had
considered that neutrinos might really be gravitons. It was, I am sure, more
endearing than useful. I haunted our local city library for books on particle
physics and cosmology; it held exactly one, and it did not help much. My parents
were not scholars and could not point the way, so I remember vividly the day my
high-school chemistry teacher mentioned that there was a \emph{scientific} library
in Copenhagen. I took the train there alone, ran my eyes along the shelves and
through the journals, and understood essentially none of it---but I came away with
something better than understanding: the certainty that all of this knowledge
existed, and was merely hard to reach.

I studied chemical engineering, won a PhD scholarship, and spent it on a doctorate
in quantum mechanics. I devoured Schr\"odinger and Boltzmann and Grafstein, the
texts on quantum thermodynamics and scattering theory. I always admired general
relativity from a distance, though I found the mathematics of its manifolds hard
to grasp and harder still to use. Afterwards I worked at the national
supercomputer centre, visualising more quantum mechanics, computational fluid
dynamics, and---less romantically---buildings and landscapes.

I found my way back toward science, in a more applied role, with the NorduGrid
project at the Niels Bohr Institute: how to gather, store, and process the vast
flood of data that CERN's soon-to-run accelerator would pour out. There I worked
with the same John Renner Hansen I had written to as a child. I became head of the
Nordic CERN Tier-1 centre and started a number of other data collaborations for
big science; I had a talk I gave often, about Tycho Brahe and Kepler---about the
long separation between those who gather the data and those who think with it. I
was at CERN several times a month. But I was \emph{directing} research
infrastructure, not doing the research, so I kept reading, and kept current, from
the side.

Research is also politics, and alongside my wish to help new knowledge into the
world I had an entrepreneurial streak. I left, stumbled onto bitcoin in 2011,
co-founded Kraken in San Francisco, and later founded Chainalysis---tracing dark
transactions instead of dark particles. Yet I spent a great many evenings and
Sundays thinking about physics: about the void, about quantum fluctuations, about
relativity. Quantum mechanics, nuclear fusion, and cosmology remained my favourite
form of procrastination.

Many of the ideas in \emph{Eternal Dawn} took shape over the last ten years, but
they finally fell together over a single sleepless night in Seoul, the morning
that followed, and the days after that---in between coding the algorithms for my
next adventure, GryAI. The framework, the simulations, and the conclusions came in
roughly a week; the polishing and the careful massaging took a couple of weeks
more. I could not have finished it this quickly without AI; my thanks to
Anthropic, and to GryAI.

\bigskip
\noindent\hfill\emph{Michael Gronager}
